\section*{第 \ref{ch:g9:how-species-change} 章——物种如何变化}

\begin{solution}{exo:g9:how-species-change:1}
\omterm{def:g9:heredity-and-traits:heredity}{遗传性状}存在变异；过量繁殖的后代受\omterm{def:g6:exploring-our-environment:components}{环境}制约，
而制约的作用并不均等——\omterm{def:g4:adapted-to-their-world:adaptation}{适应}者稍多地逃脱淘汰，
留下更多后代。其结果是，\omterm{def:g4:adapted-to-their-world:adaptation}{适应}性\omterm{def:g9:genes-and-dna:gene}{等位基因}所占比例
逐代增长——这就是\omterm{prop:g9:how-species-change:selection}{自然选择}；
长期累积则为\omterm{prop:g9:how-species-change:selection}{演化}。
\end{solution}

\begin{solution}{exo:g9:how-species-change:2}
变异：\omterm{def:g9:heredity-and-traits:heredity}{遗传}而来的浅色型与深色型。制约：有视力的
\omterm{prop:g3:sorting-living-things:groups}{鸟类}针对停在树皮上的蛾进行捕食。不均等的淘汰：在被烟尘
熏黑的树干上，浅色蛾易被看见而遭捕食，深色型的
比例随之上升；在干净的浅色树皮上，深色蛾易被看见，
浅色型比例便回升。同一机制，两个方向。
\end{solution}

\begin{solution}{exo:g9:how-species-change:3}
就\omterm{def:g8:reproduction-and-environment:population}{种群}而言，没有哪只蛾会“变深”。\omterm{prop:g9:how-species-change:selection}{自然选择}
改变的是各\omterm{def:g9:heredity-and-traits:heredity}{遗传}型式在多代间的\emph{频率}——
个体如牌般被抽出、用尽；
真正漂移的是牌组的构成。
\end{solution}

\begin{solution}{exo:g9:how-species-change:4}
努力不能\omterm{def:g9:heredity-and-traits:heredity}{遗传}——拉长的\omterm{def:g1:my-body:parts}{脖子}，好比铁匠练出的臂膀。
机制是：\omterm{def:g1:my-body:parts}{脖子}的长短本可\omterm{def:g9:heredity-and-traits:heredity}{遗传}变异；颈长的变异体
取食位置更高，在饥荒的筛选中更常存活，
留下更多携带长颈\omterm{def:g9:genes-and-dna:gene}{等位基因}的后代——
\omterm{def:g8:reproduction-and-environment:population}{种群}的\omterm{def:g1:my-body:parts}{脖子}因此变长，
却没有哪一只长颈鹿的\omterm{def:g1:my-body:parts}{脖子}因此变长。
\end{solution}

\begin{solution}{exo:g9:how-species-change:5}
选择的“引擎”按世代运转，而\omterm{def:g6:producing-our-food:fermentation}{微生物}的一代只需数小时，
一天之内便可运转数千次。若再加入一种猛烈的制约（抗生素），
罕见的抗性\omterm{prop:g9:genes-and-dna:explains}{突变}体便会把这一优势\omterm{def:g9:heredity-and-traits:heredity}{遗传}下去，
一周之内便占据了整个群体。
\end{solution}

\begin{solution}{exo:g9:how-species-change:6}
两者都是对\omterm{def:g9:heredity-and-traits:heredity}{遗传}变异施加的不均等制约——
决定谁能繁殖。区别在于：育种者有意地朝着目标选择；
\omterm{prop:g3:sorting-living-things:groups}{鸟类}并未作出任何选择——
其捕食所产生的后果完成了筛选。
\end{solution}

\begin{solution}{exo:g9:how-species-change:7}
事实之一：它只能比较\emph{现存的}变异体。
\omterm{prop:g9:how-species-change:selection}{自然选择}像一只筛子，作用于现有的牌组——它无法召来
尚未\omterm{prop:g9:genes-and-dna:explains}{突变}、未经重组的更优设计，只能从手头已有的材料中
挑出最好的；因此，肢体的设计只能是“旧部件新用”，
而非重新设计。
\end{solution}

\begin{solution}{exo:g9:how-species-change:8}
变异：喙的深浅可以\omterm{def:g9:heredity-and-traits:heredity}{遗传}（环志鸟的家系可以证明）。
制约：干旱之后，剩下的都是坚硬的\omterm{def:g2:from-seed-to-plant:seed}{种子}——
小喙的鸟啃不开剩余的食物。逐一核算：深喙的鸟得以觅食、
存活、繁殖；下一季的雏鸟平均喙更深，深喙\omterm{def:g9:genes-and-dna:gene}{等位基因}的占比
随之上升。这里并未发生分化：
只是一个\omterm{def:g8:reproduction-and-environment:population}{种群}、一次运转，且得以测量。
\end{solution}

\begin{solution}{exo:g9:how-species-change:9}
隔离：海洋——漂流者只能彼此繁殖。分化：
每个岛屿的\omterm{def:g6:exploring-our-environment:conditions}{条件}以其自身方式驱动这台“引擎”，
逐代积累各自的\omterm{def:g9:genes-and-dna:gene}{等位基因}。
检验：重新相遇的\omterm{prop:g3:sorting-living-things:groups}{鸟类}无法相互\omterm{prop:g4:animal-reproduction:where}{交配}——
按照\omterm{def:g5:biodiversity:species}{物种}的定义，这一\omterm{def:g1:plants-around-us:tree}{分支}已经完成。
\end{solution}

\begin{solution}{exo:g9:how-species-change:10}
两者施加的制约都参差不齐。未完成的疗程杀死敏感个体，
放出幸存者；这些幸存者富含抗性，而选择仍在进行之中。
农场中散布的零星剂量，每天都在庞大的\omterm{def:g6:producing-our-food:fermentation}{微生物}\omterm{def:g8:reproduction-and-environment:population}{种群}中
进行同样的部分筛选。
制约只淘汰、却不彻底，
等于为被保留下来的个体进行育种。
\end{solution}

\begin{solution}{exo:g9:how-species-change:11}
前文关于制约的章节已指出，每一代的过剩个体大多被淘汰——
或被捕食，或因饥饿与\omterm{def:g6:exploring-our-environment:components}{环境}\omterm{def:g6:exploring-our-environment:conditions}{条件}而被淘汰；关于独特性的章节则指出，
被淘汰者在\omterm{def:g9:heredity-and-traits:heredity}{遗传}上彼此不同，且逐批各异。
选择并非另外附加的东西：它不过是这样一种观察——
这些制约落在那一变异之上时，
不可能平等地落下；
两章内容相乘，便是\omterm{prop:g9:how-species-change:selection}{自然选择}。
\end{solution}

\begin{solution}{exo:g9:how-species-change:12}
选择消耗变异；洗牌与\omterm{prop:g9:genes-and-dna:explains}{突变}则为其补充储备。克隆\omterm{def:g8:reproduction-and-environment:population}{种群}只给
筛子提供一副牌：\omterm{def:g6:exploring-our-environment:components}{环境}一旦改变，选择便无可提升的对象，
葡萄园要么整体兴衰，因为它本身只是一株\omterm{def:g1:plants-around-us:plant}{植物}。
若没有新\omterm{def:g9:genes-and-dna:gene}{等位基因}的不断涌入，
也没有新组合的不断发牌，
引擎便会停摆：变异不是机制的副产品，
而是它的燃料。
\end{solution}

\begin{solution}{pb:g9:how-species-change:1}
\textbf{1.} \omterm{def:g6:how-plants-spread:dispersal}{扩散}一章指出：只有远距离的“运输方式”——翅膀、漂流、
风暴吹送者——才能抵达大洋岛屿；
岛上的\omterm{def:g5:biodiversity:species}{物种}名单，
即这些运输方式的“舱单”。

\textbf{2.} 群岛上的雀类源自大陆的漂流者——
细节上的共有特征意味着共同的祖先，
其最近的亲缘者
即指明了来源\omterm{def:g8:reproduction-and-environment:population}{种群}。

\textbf{3.} \omterm{def:g4:adapted-to-their-world:adaptation}{适应}一章讲的是“合适”——
每种喙都回应其所在岛屿的食物问题。本章补充
的是\emph{形成过程}：漂流者的变异、
岛上的选择\omterm{def:g6:exploring-our-environment:conditions}{条件}、不均等的采集、世代相传的记账——
“合适”的制造过程。此前我们只能看到“合适”
本身，
（\cref{rem:g4:adapted-to-their-world:how}所承诺的
答案）。

\textbf{4.} 隔离：被海洋阻隔的\omterm{def:g8:reproduction-and-environment:population}{种群}从一开始就各自繁殖，
每座岛屿都是一个密封的发动机舱，
——\omterm{def:g5:biodiversity:species}{物种}形成的第一要素，
由地理整批提供。

\textbf{5.} 选择\omterm{def:g6:exploring-our-environment:conditions}{条件}到来（仅硬\omterm{def:g2:from-seed-to-plant:seed}{种子}可食）；
喙深的\omterm{def:g9:heredity-and-traits:heredity}{遗传}变异被不均等地采集
（小喙者饿死），存活者的繁殖
把这一变化写入下一代的平均值——
发动机的一次运转，卡尺有据可查。

\textbf{6.} 其“方向”不过是当下\omterm{def:g6:exploring-our-environment:conditions}{条件}所致：
\omterm{prop:g9:how-species-change:selection}{自然选择}跟随\omterm{def:g6:exploring-our-environment:conditions}{条件}而变，
\omterm{def:g6:exploring-our-environment:conditions}{条件}一旦逆转，漂移也随之逆转——没有终点，
只有逐季而来的后果。

\textbf{7.} \Cref{prop:g9:how-species-change:speciation}
进行到中途：分化已推进到
杂交开始受阻的程度——一道闭合中的分叉，
就在笔记的注视下逐渐合拢。

\textbf{8.} \omterm{def:g5:biodiversity:species}{物种}形成的“滴答”是数以千计的世代——
只有快速更替的生命（病房里的\omterm{def:g6:producing-our-food:fermentation}{微生物}）
才能被完整观察。对于雀类，
这一方法转而以旁证的形式出现：格局、\omterm{prop:g6:classification-kinship:kinship}{亲缘关系}、
活体的“滴答”与半合拢的分叉，
共同指向任何单一\omterm{prop:g2:born-grow-die:lifespan}{寿命}都无法跨越的过程。

\textbf{9.} （1）变异：漂流而来的奠基者，经重组并发生\omterm{prop:g9:genes-and-dna:explains}{突变}。
（2）选择\omterm{def:g6:exploring-our-environment:conditions}{条件}：各岛自有的食物与危险。
（3）记账：各岛适宜的\omterm{def:g9:genes-and-dna:gene}{等位基因}逐岛累积。
（4）隔离：其间的海洋使各本账目彼此分开。
（5）时间尺度：数千代——
而分叉此时已然完成。

\textbf{10.} 只有其奠基者所携带的，加上此后自身的\omterm{prop:g9:genes-and-dna:explains}{突变}：
每座岛屿的“牌组”，都是一名漂流者的样本，
\omterm{prop:g9:how-species-change:selection}{自然选择}只能从中挑出\emph{现有}的最优喙形——
牌组更丰富、或\omterm{prop:g9:genes-and-dna:explains}{突变}更走运的岛屿，
得到的答案就更接近。

\textbf{11.} 世代时间：病房里是数小时，
岛上则是数年——相同的三项事实、相同的记账，
滴答之差达数千倍。

\textbf{12.} 两者皆是决定谁来繁殖的选择\omterm{def:g6:exploring-our-environment:conditions}{条件}：
饲养者的有意之手，与干旱的无意识选择
只留下硬\omterm{def:g2:from-seed-to-plant:seed}{种子}，对\omterm{def:g9:heredity-and-traits:heredity}{遗传}变异做的是同样的记账——
这正是鸽舍里数百年间造成的变化，
足以证明自然的选择
在数百万年间能写下怎样的结果。
\end{solution}
